\documentclass[10pt,a4paper]{amsart}
\usepackage[margin=19mm]{geometry}
\usepackage{amsmath,amssymb,mathtools}
\usepackage[scr=rsfs,cal=euler]{mathalfa}
\usepackage{xfrac}
\usepackage[unicode,hidelinks,bookmarks=false]{hyperref}
\newcommand{\ie}{i.e.\ }
\newcommand{\cf}{cf.\ }
\newcommand{\resp}{respectively\ }
\newcommand{\Subsection}{\subsection}
\providecommand{\nobreakdash}{\nobreak\hbox{-}}
\setlength{\emergencystretch}{2em}
\multlinegap0pt
\usepackage{bm}
\usepackage{bbm}
\usepackage{dsfont}
\usepackage{xspace}
\usepackage{cancel}
\usepackage{mathtools}
\usepackage{amsthm}
\makeatletter
\@ifundefined{theorem}{\newtheorem{theorem}{Theorem}}{}
\@ifundefined{proposition}{\newtheorem{proposition}[theorem]{Proposition}}{}
\makeatother
\title[On quotients of ideals of weighted holomorphic mappings]{On quotients of ideals of weighted holomorphic mappings}
\author{Belacel Amar, Bougoutaia Amar, Rueda Pilar}
\date{Source: arXiv:2505.19593v1 (CC BY 4.0)}
\begin{document}
\maketitle

\noindent\emph{Source and licence.} On quotients of ideals of weighted holomorphic mappings. Version arXiv:2505.19593v1 (CC BY 4.0), distributed under the Creative Commons Attribution 4.0 International licence, as recorded in the arXiv licence metadata for this exact version and shown on the abstract page https://arxiv.org/abs/2505.19593v1. Full rights notice: licenses/arxiv-2505.19593-CC-BY-4.0.txt.\par\smallskip

\noindent\textit{Editorial scope.} Counted unit: the Proposition whose source label is \texttt{coincidence} in the article (source0.tex lines 286--306, source label \texttt{coincidence}), delivered together with the article's own complete original proof of it (source0.tex lines 308--362, one \texttt{proof} environment of 55 lines). No mathematics is written, repaired or re-derived: statement and proof are the article's own text line for line. Required context retained uncounted: none. Every definition, display and label of those lines is retained unchanged. Omitted from this package and not redistributed: 1--285, 307--307, 363--883. Nothing inside a retained excerpt is omitted or altered: every retained body is delivered exactly as the article's lines are written, so no omission receipt of any kind accompanies this package. Citations: the retained excerpts carry no citation command at all, so no bibliography entry of the article is transcribed and no reference is redistributed. Reference adaptation: none was needed and none was made; every reference inside the delivered unit resolves inside it, so no unresolved reference or citation is produced. Printed slips kept literal: no printed word, symbol, hypothesis or number of the counted statement or of its proof was altered. Layout and class: the article's own class is \texttt{amsart} with packages including eurosym, amscd, amsmath, amssymb, enumerate, amsfonts, mathrsfs, txfonts, comment, hyperref, xcolor, tikz, pgfplots, pgf; the wrapper is the lane's pinned \texttt{amsart} editorial wrapper at 10pt on A4, which restates the theorem-like environments the retained text needs and adds the article's own macro definitions that the retained text needs (none), with extra wrapper packages (bm, bbm, dsfont, xspace, cancel, mathtools). The delivered numbering is the unit's own. Assets: no retained span contains a figure environment, a graphics-inclusion command, a PDF-page inclusion, a table or a TikZ picture; no image file is copied into this package, the package declares no asset dependency and no image file is redistributed with it. Licence: version arXiv:2505.19593v1 of this article is distributed under the Creative Commons Attribution 4.0 International licence, as recorded in the arXiv licence metadata for this exact version and shown on the abstract page https://arxiv.org/abs/2505.19593v1. A full rights notice is delivered at licenses/arxiv-2505.19593-CC-BY-4.0.txt. Characters: every retained excerpt is pure ASCII, so the delivered engine, XeLaTeX with its default Latin Modern text font, reports no missing character.
\begin{proposition}\label{coincidence}
\label{3}Let $\left[ \mathcal{A},\left\Vert .\right\Vert _{\mathcal{A}}%
\right] $ be a Banach operator ideal and let $\left[ \mathcal{I}^{\mathcal{H}%
_{v}^{\infty }},\left\Vert .\right\Vert _{\mathcal{I}^{\mathcal{H}%
_{v}^{\infty }}}\right] $ be a normed weighted holomorphic ideal. Then,%
\begin{equation*}
\left[ \mathcal{A}^{-1}\mathcal{\circ I}^{\mathcal{H}_{v}^{\infty
}},\left\Vert .\right\Vert _{\mathcal{A}^{-1}\mathcal{\circ I}^{\mathcal{H}%
_{v}^{\infty }}}\right] \leq \left[ \mathcal{H}_{v}^{\infty },\left\Vert
.\right\Vert _{v}\right] .
\end{equation*}

Furthermore%
\begin{equation*}
\left[ \mathcal{A}^{-1}\mathcal{\circ I}^{\mathcal{H}_{v}^{\infty
}},\left\Vert .\right\Vert _{\mathcal{A}^{-1}\mathcal{\circ I}^{\mathcal{H}%
_{v}^{\infty }}}\right] =\left[ \mathcal{H}_{v}^{\infty },\left\Vert
.\right\Vert _{v}\right]
\end{equation*}%
if $\mathcal{I}^{\mathcal{H}_{v}^{\infty }}$ has the LP in $\mathcal{A}$.
\end{proposition}

\begin{proof}
Let $f\in \mathcal{A}^{-1}\mathcal{\circ I}^{\mathcal{H}_{v}^{\infty
}}\left( U,F\right) $, we have that $\left\Vert f\right\Vert _{v}\leq
\left\Vert f\right\Vert _{\mathcal{A}^{-1}\mathcal{\circ I}^{\mathcal{H}%
_{v}^{\infty }}}$. Indeed, notice that $f\in \mathcal{H}_{v}^{\infty }\left(
U,F\right) $ and $A\circ f\in \mathcal{I}^{\mathcal{H}_{v}^{\infty }}\left(
U,G\right) $ for all $A\in \mathcal{A}\left( F,G\right) $ and all 
complex Banach space $G$. For any $x\in U$ consider a functional $y^{\ast }_x$ in the closed unit ball of the dual space $F^{\ast }$ of $F$ so that $\left\Vert
f\left( x\right) \right\Vert =\left\vert y^{\ast }_x\left( f\left( x\right)
\right) \right\vert $. Taking into account that $\left[ \mathcal{A}%
,\left\Vert .\right\Vert _{\mathcal{A}}\right] $ is an operator ideal,
we have that  $y^{\ast }_x$ belongs to $\mathcal{A}\left( F,\mathbb{C}\right) $
with $\left\Vert y^{\ast }_x\right\Vert _{\mathcal{A}}=\left\Vert
y^{\ast }_x\right\Vert \leq 1$. Then, $y^*_x\circ f\in \mathcal{ I}^{\mathcal{H}_{v}^{\infty
}}$ and $\|y_x^{\ast}\circ f\|_{\mathcal I^{\mathcal H_v^\infty}}\leq \|f\|_{\mathcal A^{-1}\circ \mathcal I^{\mathcal H_v^\infty}}$. Hence, 

$$
v(x)\|f(x)\|=v(x)|y_x^*(f(x))|\leq \|y_x^*\circ f\|_v\leq \|y_x^*\circ f\|_{\mathcal I^{\mathcal H_v^\infty}}\leq \|f\|_{\mathcal A^{-1}\circ \mathcal I^{\mathcal H_v^\infty}}
$$
and taking the supremum over all $x\in U$, we conclude that $\left\Vert
f\right\Vert _{v}\leq \left\Vert f\right\Vert _{\mathcal{A}^{-1}\mathcal{%
\circ I}^{\mathcal{H}_{v}^{\infty }}}.$

Let us now assume that $\mathcal{I}^{\mathcal{H}_{v}^{\infty }}$ has the 
\textit{LP} in the operator ideal $\mathcal{A}$. Let $f\in \mathcal{H}%
_{v}^{\infty }\left( U,F\right) $ and let $A\in \mathcal{A}\left( F,G\right) 
$, for a complex Banach space $G$. It is not difficult to see that $A\circ
f\in \mathcal{H}_{v}^{\infty }\left( U,G\right) $. Consider the linearizations $T_f\in \mathcal{L}(\mathcal{G}_{v}^{\infty }(U),F)$ and $T_{A\circ
f}\in \mathcal{L}(\mathcal{G}_{v}^{\infty }(U),G)$. We know that  $\left\Vert
T_{A\circ f}\right\Vert =\left\Vert A\circ f\right\Vert _{v}$ and $\left\Vert T_{f}\right\Vert =\left\Vert f\right\Vert _{v}$. Since
\begin{equation*}
T_{A\circ f}\circ \Delta _{v}=A\circ f=A\circ T_{f}\circ \Delta _{v},
\end{equation*}
by the uniqueness of the linearization, $T_{A\circ f}=A\circ T_{f}$. Thus, by the ideal property of $\mathcal{A}$,
we have $T_{A\circ f}\in \mathcal{A}(\mathcal{G}_{v}^{\infty }(U),G)$ with%
\begin{equation*}
\left\Vert T_{A\circ f}\right\Vert _{\mathcal{A}}=\left\Vert A\circ
T_{f}\right\Vert _{\mathcal{A}}\leq \left\Vert A\right\Vert _{\mathcal{A}%
}\left\Vert T_{f}\right\Vert \leq \left\Vert A\right\Vert _{\mathcal{A}%
}\left\Vert f\right\Vert _{v}.
\end{equation*}%
Since $\mathcal{I}^{\mathcal{H}_{v}^{\infty }}$ has the \textit{LP} in $%
\mathcal{A}$  we have $A\circ f\in \mathcal{I}^{\mathcal{H}_{v}^{\infty }}\left(
U,G\right) $ with $\left\Vert A\circ f\right\Vert _{\mathcal{I}^{\mathcal{H}%
_{v}^{\infty }}}=\left\Vert T_{A\circ f}\right\Vert _{\mathcal{A}}$. Therefore,  $f\in 
\mathcal{A}^{-1}\mathcal{\circ I}^{\mathcal{H}_{v}^{\infty }}\left(
U,F\right) $ and 
\begin{eqnarray*}
\left\Vert f\right\Vert _{\mathcal{A}^{-1}\mathcal{\circ I%
}^{\mathcal{H}_{v}^{\infty }}}&=&\sup\{ \|A\circ f\|_{\mathcal I^{\mathcal H_v^\infty}} : A\in \mathcal A, \|A\|_{\mathcal A}\leq 1\}\\
&=&  \sup\{ \|T_{A\circ f}\|_{\mathcal A} : A\in \mathcal A, \|A\|_{\mathcal A}\leq 1\}\\
&=&\sup\{ \|A\circ T_f\|_{\mathcal A} : A\in \mathcal A, \|A\|_{\mathcal A}\leq 1\}\leq \|T_f\|= \left\Vert f\right\Vert _{v}
 \end{eqnarray*}
  as desired.
\end{proof}

\end{document}
